\documentclass[9pt]{slides} % Use option 'handout' for handouts
                                      % Use option 'notes=only' for notes

\title
{Presentation title}

\subtitle
{Secondary title}

\author
{Ravi Rajani}

\institute
{Department of Mathematics}

\titlegraphic
{\includegraphics[height=0.8cm]{aselogo_en.pdf}}

\date{\today}

\subject{Subject of talk}

\begin{document}
\setbeamercolor{background canvas}{bg=Core}
\begin{frame}
  \titlepage
\end{frame}
\setbeamercolor{background canvas}{bg=white}

\section{First Section}

\begin{frame}{Outline}
  \tableofcontents[currentsection]
\end{frame}

\subsection{Subsection}

\begin{frame}{Colours}
\begin{theorem}[Central Limit Theorem]
You know it. \note{explain if not}
\end{theorem}
\begin{proof}[Proof Sketch]
Just use characteristic functions.
\end{proof}
\begin{example}
Example environment
\end{example}
\begin{minipage}{.5\textwidth}
  \begin{exampleblock}{Custom Example}
    Example body.
    \end{exampleblock}
\end{minipage}
\begin{alertblock}{Careful}
  body text.
\end{alertblock}
\begin{plainblock}{Plain Title}
This is a plain block.
\end{plainblock}
\end{frame}

\section{Second Section}

\begin{frame}[fragile]{Fragile Option and Text Positioning}
    Remember:
    \begin{itemize}[<+->]
      \item Use the \texttt{fragile} frame option for \textcolor<2->{OUMidBlue}{verbatim} environments.
      \item Use \colorbox{OUYellow}{background templates} and \alert{texpos} as follows:
            \begin{verbatim}
\setbeamertemplate{background}{\import{subfolder}{file.pdf}}
\begin{frame}
  % The textblock* environment does absolute positioning
  % First param is width of box, and second is
  % coordinates of left corner in ROUND brackets!
  \begin{textblock*}{7cm}(7cm,1.5cm)
    ...
    
  \end{textblock*}
\end{frame}% Reset the background template after:
\setbeamertemplate{background}{}
\end{verbatim}
    \end{itemize}
\end{frame}

\setbeamercolor{background canvas}{bg=Core}
\begin{frame}
  \usebeamertemplate{end page}
\end{frame}

\end{document}
